% ============================================================
%  PREAMBLE
% ============================================================
\documentclass[12pt, a4paper]{article}   % font size, paper size

% Encoding & language
\usepackage[utf8]{inputenc}
\usepackage[T1]{fontenc}
\usepackage[english]{babel}

% Page layout
\usepackage[margin=2.5cm]{geometry}

% Math
\usepackage{amsmath}    % extended math environments
\usepackage{amssymb}    % extra symbols (ℝ, ∞, etc.)

% Graphics & figures
\usepackage{graphicx}   % \includegraphics
\usepackage{float}      % [H] placement option

% Tables
\usepackage{booktabs}   % \toprule, \midrule, \bottomrule

% Hyperlinks (load last)
\usepackage[hidelinks]{hyperref}

% ---- Document metadata ----
\title{Title}
\author{Author Name}
\date{\today}           % or a fixed date: \date{15 May 2026}


% ============================================================
%  BODY
% ============================================================
\begin{document}

\maketitle              % renders title, author, date
\tableofcontents        % auto-generated from \section{} headings
\newpage


% ---- Section ----
\section{Introduction}

Body text goes here. A blank line in the source starts a new paragraph.
Inline math: $E = mc^2$.

% ---- Subsection ----
\subsection{Background}

\subsubsection{Detail}

% ---- Unnumbered section (won't appear in ToC by default) ----
\section*{Unnumbered Section}


% ---- Lists ----
\begin{itemize}
    \item First bullet
    \item Second bullet
        \begin{itemize}
            \item Nested bullet
        \end{itemize}
\end{itemize}

\begin{enumerate}
    \item First item
    \item Second item
\end{enumerate}


% ---- Display math ----
\begin{equation}
    \int_a^b f(x)\,dx = F(b) - F(a)
    \label{eq:fundamental}
\end{equation}

See Equation~\ref{eq:fundamental}.   % cross-reference

% Multi-line aligned equations
\begin{align}
    y &= mx + b \\
    z &= ax^2 + bx + c
\end{align}


% ---- Figure ----
\begin{figure}[H]
    \centering
    \includegraphics[width=0.7\textwidth]{/Users/johannes/Downloads/UWC_RBC_zweizeilig_unten_CMYK.png}   % no extension needed
    \caption{Caption text.}
    \label{fig:example}
\end{figure}

See Figure~\ref{fig:example}.


% ---- Table ----
\begin{table}[H]
    \centering
    \caption{Caption text.}
    \label{tab:example}
    \begin{tabular}{lcc}
        \toprule
        Column A    & Column B  & Column C  \\
        \midrule
        Row 1       & value     & value     \\
        Row 2       & value     & value     \\
        \bottomrule
    \end{tabular}
\end{table}

See Table~\ref{tab:example}.


% ---- Bibliography ----
% Option A: manual
\begin{thebibliography}{9}
    \bibitem{key1} Author, \textit{Title}, Publisher, Year.
\end{thebibliography}

% Option B: BibLaTeX (comment out Option A and uncomment below)
% \usepackage{biblatex}          % add to preamble
% \addbibresource{refs.bib}      % add to preamble
% \printbibliography             % replaces \begin{thebibliography}

\end{document}